% =============================================================================
% Paper 1: Ternary Edge-RV: Full-Stack Ternary NPU on RISC-V Linux
% Template estilo SBCCI / LASCAS (double-column, 4-6 pages)
% Última atualização: 20/08/2026
% =============================================================================

\documentclass[10pt,conference]{IEEEtran}

% PACOTES
\usepackage[utf8]{inputenc}
\usepackage[T1]{fontenc}
\usepackage{graphicx}
\usepackage{amsmath,amssymb}
\usepackage{cite}
\usepackage{url}
\usepackage{booktabs}
\usepackage{multirow}
\usepackage{algorithmic}
\usepackage{algorithm}

% META-INFORMAÇÃO
\hyphenation{ter-na-ry no-nu-cle-ar mul-ti-pli-er-less}

\title{Ternary Edge-RV: A Multiplierless Ternary Neural Network Accelerator for RISC-V Linux Systems}

\author{
    \ IEEEauthorblockN{Arthur Oliveira Gomes$^{1}$, Gildo Alves de Lima Junior$^{2}$, \\
    Gustavo Alexandre dos Santos$^{3}$, Gilvan Alves Pastor Junior$^{4}$}
    \ IEEEauthorblockA{Departamento de ... , Universidade Federal de ... \\
    \{arthur.gomes, gildo.junior, gustavo.santos, gilvan.pastor\}@email.br}
}

\begin{document}

\maketitle

% =============================================================================
% ABSTRACT
% =============================================================================
\begin{abstract}
Ternary Neural Networks (TNNs) constrain weights to $\{-1, 0, 1\}$, reducing ternary products to conditional additions and subtractions for edge inference. However, most existing TNN accelerators require either standalone FPGA logic without a host OS or tightly coupled custom ASICs. We present \textbf{Ternary Edge-RV}, a prototype full-stack system targeting integration of a ternary NPU as a Wishbone peripheral in a LiteX-generated RV32IMA SoC running Linux. The source-level architecture targets multiplierless ternary MAC processing, multi-layer sequencing, and Wishbone Master DMA for direct memory access. Host-side validation exists for selected models and software interfaces, but trained-model provenance, physical synthesis, FPGA deployment, accuracy, latency, throughput, CPU utilization, and DSP usage remain TBD until the corresponding evidence is available.
\end{abstract}

\begin{IEEEkeywords}
Ternary neural networks, RISC-V, FPGA, multiplierless, NPU, edge computing, LiteX.
\end{IEEEkeywords}

% =============================================================================
% INTRODUCTION
% =============================================================================
\section{Introduction}
\label{sec:intro}

% TODO: Escrever introdução (equipe toda)
% Contexto: Edge AI, TinyML, restrições de energia em dispositivos IoT
% Problema: Multiplicadores em FPGA são caros (DSP slices limitados)
% Solução proposta: TNN + RISC-V Linux = sistema completo open-source
% Contribuições: (listar 3-4 contribuições principais)

Edge artificial intelligence (AI) inference on resource-constrained devices demands architectures that maximize computational efficiency within strict area and power budgets. While CPUs provide flexibility, their general-purpose nature incurs high energy overhead for the repetitive multiply-accumulate (MAC) operations dominant in neural networks. Conversely, dedicated neural processing units (NPUs) achieve superior efficiency but often lack the software ecosystem of a general-purpose OS.

Ternary Neural Networks (TNNs) \cite{ternarynet} offer a unique opportunity: by restricting weights to the set $\{-1, 0, 1\},$ multiplications reduce to conditional additions and subtractions, completely eliminating the need for hardware multipliers. This makes TNNs ideally suited for FPGA deployment where DSP slices are scarce.

This paper presents \textbf{Ternary Edge-RV}, a prototype full-stack system targeting integration of a multiplierless ternary NPU into a Linux-capable RISC-V system-on-chip (SoC). Our contributions are:

\begin{enumerate}
    \item \textbf{Source-level multiplierless NPU target} organized around up to 64 parallel MAC units using adders, subtracters, and multiplexers; parallelism and DSP utilization remain subject to RTL and synthesis validation.
    \item \textbf{Wishbone Master DMA} intended to enable autonomous memory access and reduce CPU involvement during inference.
    \item \textbf{Hardware layer sequencer} intended to chain multiple layers without CPU intervention.
    \item \textbf{Full-stack integration path} from hardware (LiteX/VexRiscv) through Linux kernel driver to user-space AI application.
\end{enumerate}

% =============================================================================
% BACKGROUND & RELATED WORK
% =============================================================================
\section{Background and Related Work}
\label{sec:background}

% TODO: Arthur + Gustavo (active AI pipeline maintenance) + Gilvan (historical QAT, packing, and Golden Model contribution)
% - Ternary Neural Networks (quantização, STE, benefícios)
% - Multiplierless MAC architecture
% - Trabalhos relacionados: FINN (Xilinx), Larq, outras NPUs RISC-V

\subsection{Ternary Neural Networks}
Ternary Neural Networks \cite{ternarynet, trainedternary} constrain weights to $\{-1, 0, 1\}$, typically using a threshold-based quantization during training:

\begin{equation}
w_q = \begin{cases}
+1 & \text{if } w > \Delta \\
0  & \text{if } |w| \leq \Delta \\
-1 & \text{if } w < -\Delta
\end{cases}
\end{equation}

During backpropagation, the Straight-Through Estimator (STE) \cite{binaryconnect} approximates the gradient of the non-differentiable quantization function.

\subsection{Existing FPGA Accelerators}
Previous work such as FINN \cite{finn} and Larq Compute Engine \cite{larq} demonstrated binary and ternary networks on FPGAs, but either require specialized Xilinx tools (FINN) or lack integration with a full OS. Our approach targets generic Wishbone-based RISC-V SoCs, maximizing portability.

% =============================================================================
% SYSTEM ARCHITECTURE
% =============================================================================
\section{System Architecture}
\label{sec:arch}

% TODO: Arthur (hardware) + Gildo (OS/DT) + Gustavo (driver, AI pipeline/export, Golden Model regression, benchmarks, and results) + Gilvan (historical QAT/packing/Golden Model contribution)
% Esta é a seção mais importante: cada um escreve sua parte

\subsection{Hardware: Multiplierless NPU}
\label{sec:hardware}

% --- Arthur escreve esta subseção ---
The source-level design integrates the NPU as a Wishbone B4 slave peripheral in a LiteX-generated SoC featuring a VexRiscv core configured as RV32IMA (32-bit with atomic instructions and MMU). Figure~\ref{fig:arch} shows the planned top-level system diagram.

\begin{figure}[htbp]
\centering
% TODO: Inserir diagrama de blocos do sistema
% \includegraphics[width=\columnwidth]{figures/system_architecture.pdf}
\vspace{3cm}
\caption{Top-level system architecture showing the RISC-V SoC, NPU, and DMA data flow.}\label{fig:arch}
\end{figure}

\subsubsection*{The Ternary MAC}
The core computational unit is the \textit{ternary MAC}, which replaces a standard $8$-bit $\times$ $8$-bit multiplier with a conditional adder:

\begin{equation}
P(i) = 
\begin{cases}
+act[i] & \text{if } w[i] = 01\;(+1) \\
0       & \text{if } w[i] = 00\;(0) \\
-act[i] & \text{if } w[i] = 11\;(-1)
\end{cases}
\end{equation}

\subsubsection*{64-MAC Array}
The source-level NPU v2 architecture targets up to 64 ternary MAC units operating in parallel. The intended datapath fetches four 32-bit words (128 bits total) from the internal weight buffer and unpacks 64 $\times$ 2-bit weights for the MAC array. The active top-level currently uses a procedural accumulation path, so the 64-MAC parallelism and pipelined adder-tree realization are architectural targets rather than validated implementation results. RTL and synthesis evidence are required before reporting parallel throughput, a 50 MHz operating point, or DSP utilization.

\subsubsection*{Wishbone Master DMA}
The proposed source-level accelerator includes a Wishbone B4 Master controller intended to perform burst reads directly from system RAM. The driver is intended to program source and destination addresses via MMIO registers, after which the NPU should fetch weights and activations, process the configured layers, and write the result back to RAM before raising an interrupt. End-to-end DMA and interrupt behavior remain pending physical validation.

\subsubsection*{Layer Sequencer}
A dedicated finite state machine (FSM) is intended to iterate through the three dense layers (784$\to$1024, 1024$\to$512, 512$\to$256) automatically. The intended sequencer reconfigures the DMA controller between layers without CPU intervention, subject to RTL and system validation.

\subsection{Memory-Mapped Register Map}
\label{sec:regmap}

% --- Arthur escreve esta subseção ---
% NOTE: 0x80000000 is the current LiteX NPU MMIO base. The earlier 0x40000000 map is historical/target-specific and must be reconciled with the generated SoC map before physical deployment.

Table~\ref{tab:regmap} lists the memory-mapped registers accessible from Linux via the kernel driver.

\begin{table}[htbp]
\centering
\caption{NPU register map (current LiteX base address 0x8000\_0000).}
\label{tab:regmap}
\small
\begin{tabular}{clll}
\toprule
\textbf{Offset} & \textbf{Register} & \textbf{R/W} & \textbf{Description} \\
\midrule
0x00 & STATUS & RO & Busy, IRQ pending, zero counter \\
0x04 & CONTROL & WO & Start, clear IRQ \\
0x08 & DMA\_SRC\_ADDR & R/W & Source physical address \\
0x0C & DMA\_DST\_ADDR & R/W & Destination physical address \\
0x10 & DMA\_SIZE & R/W & Total MAC operations \\
0x14 & WEIGHT\_CFG & R/W & Weight row config \\
0x18 & ACT\_CFG & R/W & Activation config \\
0x1C & RESULT & RO & Final result \\
\bottomrule
\end{tabular}
\end{table}

\subsection{Operating System and Device Tree}
\label{sec:os}

% --- Gildo escreve esta subseção ---

The target SoC is configured for Linux 6.x built with Buildroot. The kernel configuration targets RV32IMA with high-resolution timers enabled. The NPU is declared in the Device Tree:

\begin{verbatim}
npu_ternaria: npu@80000000 {
    compatible = "ternaryedge,npu-ternaria";
    reg = <0x80000000 0x00001000>;
    interrupts = <0x0a>;
    interrupt-parent = <&plic>;
};
\end{verbatim}

\subsection{Kernel Driver}
\label{sec:driver}

% --- Gustavo escreve esta subseção; active owner: driver, cross-compilation, and MMIO validation ---

The kernel driver implements a Platform Driver that matches against the Device Tree node. Key features include:

\begin{itemize}
    \item \textbf{MMIO mapping:} Registers accessed via \texttt{devm\_ioremap\_resource()}.
    \item \textbf{DMA buffer:} Coherent DMA memory allocated via \texttt{dma\_alloc\_coherent()} for zero-copy weight transfer.
    \item \textbf{Interrupt handling:} \texttt{devm\_request\_irq()} with a wait queue puts the user process to sleep during NPU compute.
    \item \textbf{Character device:} \texttt{/dev/npu\_ternaria} exposes \texttt{open()}, \texttt{mmap()}, \texttt{ioctl()}, and \texttt{release()}.
\end{itemize}

\subsection{AI Pipeline: From Training to Deployment}
\label{sec:ai}

% --- Gustavo escreve esta subseção; active owner: AI pipeline maintenance and export/weights.h ---
% Gilvan's historical QAT, packing, and Golden Model contribution remains credited.

\subsubsection*{Quantization-Aware Training}
The repository contains a QAT training path for a 3-layer MLP on MNIST (784$\to$1024$\to$512$\to$256$\to$10) using Larq \cite{larq}. The current status of the model and exported header requires validation against the deployment contracts:

\begin{itemize}
    \item Weights: ternary ($\{-1,0,1\}$, 2-bit encoding)
    \item Activations: INT8
    \item Regularization: L1 sparsity penalty configured as a training objective; achieved sparsity is TBD
    \item Accuracy: TBD until the trained model is validated on the MNIST test set
\end{itemize}

The provenance of the real trained model and the generated \texttt{weights.h} header requires validation. The checked-in FP32 symbols currently use fallback $0.01$/$0.1$ values rather than validated trained-model output. The export format, inter-layer transforms, and their agreement with the HAL and RTL data contracts are also pending validation; no accuracy or deployment result is claimed here.

\subsubsection*{Weight Packing}
Each 32-bit word is intended to pack 16 ternary weights in little-endian order (weight[0] maps to bits $[1:0]$). The encoding target is: $+1 = 01_2$, $0 = 00_2$, $-1 = 11_2$. The exported header and inter-layer transforms require validation. This historical packing work is retained as part of Gilvan's contribution, while Gustavo owns active export and regression maintenance.

\subsubsection*{Golden Model}
A C++ simulation of the NPU (\texttt{npu\_sim.cpp}) is retained as the host-side Golden Model for bit-accurate regression. Gilvan's historical Golden Model contribution remains credited; Gustavo owns the active regression and validation work.

% =============================================================================
% EXPERIMENTAL RESULTS
% =============================================================================
\section{Experimental Results}
\label{sec:results}

% TODO: Gustavo (cross-compilation, physical benchmarks, and results) + Arthur (synthesis)
% Tabela de síntese (LUTs, FFs, DSPs: all TBD until synthesis)
% Gráfico tempo inferência CPU vs NPU (all values TBD until physical benchmark)
% Tabela comparativa com trabalhos relacionados

Current host-side evidence includes C++ v1 validation with 8/8 tests passing, C++ v2 validation with 21/21 tests passing, Python checks with 5/5 tests passing, and a passing IOCTL ABI check. The Verilog testbench is unavailable in the current shell, and no FPGA end-to-end result is reported. Physical synthesis, deployment, accuracy, and performance measurements remain pending.

\begin{table}[htbp]
\centering
\caption{Planned synthesis results; all physical values remain TBD.}
\label{tab:synth}
\small
\begin{tabular}{lrr}
\toprule
\textbf{Resource} & \textbf{Used} & \textbf{Available} \\
\midrule
LUTs & TBD & TBD \\
FFs  & TBD & TBD \\
DSPs & TBD & TBD \\
BRAM (36Kb) & TBD & TBD \\
\midrule
Max. Freq. & TBD & TBD \\
\bottomrule
\end{tabular}
\end{table}

\begin{table}[htbp]
\centering
\caption{Planned inference time comparison: CPU (RISC-V software) vs NPU (hardware).}
\label{tab:bench}
\small
\begin{tabular}{lrrr}
\toprule
\textbf{Layer} & \textbf{CPU (ms)} & \textbf{NPU (ms)} & \textbf{Speedup} \\
\midrule
784 $\to$ 1024 & TBD & TBD & TBD \\
1024 $\to$ 512 & TBD & TBD & TBD \\
512 $\to$ 256  & TBD & TBD & TBD \\
\midrule
Total & TBD & TBD & TBD \\
\bottomrule
\end{tabular}
\end{table}

% =============================================================================
% CONCLUSION
% =============================================================================
\section{Conclusion}
\label{sec:conclusion}

% TODO: Gustavo owns pending benchmarks and results; all authors review the final status and conclusion.
We presented the prototype and source-level integration plan for Ternary Edge-RV, a ternary neural network accelerator targeting a RISC-V Linux system. Host-side validation exists for selected C++ models, Python checks, and the IOCTL ABI, while the real trained-model/header provenance and inter-layer transforms still require validation. Physical synthesis, FPGA deployment, accuracy, resource use, frequency, latency, throughput, CPU utilization, energy, and DSP utilization remain pending gates, so no physical performance result is claimed. The prototype provides a basis for the planned end-to-end evaluation once these gates are completed.

% =============================================================================
% REFERENCES
% =============================================================================
\begin{thebibliography}{00}

\bibitem{ternarynet} F. Li, B. Zhang, and B. Liu, ``Ternary weight networks,'' \textit{arXiv preprint arXiv:1605.04711}, 2016.

\bibitem{trainedternary} D. Wan, J. Shen, \textit{et al.}, ``Trained ternary quantization,'' in \textit{Proc. ICLR}, 2017.

\bibitem{binaryconnect} M. Courbariaux, Y. Bengio, and J.-P. David, ``BinaryConnect: Training deep neural networks with binary weights during propagations,'' in \textit{Proc. NeurIPS}, 2015.

\bibitem{finn} Y. Umuroglu, N. J. Fraser, \textit{et al.}, ``FINN: A framework for fast, scalable binarized neural network inference,'' in \textit{Proc. FPGA}, 2017.

\bibitem{larq} L. Geiger and P. Team, ``Larq: An open-source library for training binarized neural networks,'' \textit{Journal of Open Source Software}, vol. 5, no. 46, 2020.

\end{thebibliography}

\end{document}
